\documentclass[12pt,a4paper]{article}

% ── Packages ──────────────────────────────────────────────────────────────────
\usepackage[margin=1in]{geometry}
\usepackage{amsmath,amssymb,mathtools}
\usepackage{booktabs}
\usepackage{enumitem}
\usepackage{hyperref}
\usepackage{titlesec}
\usepackage{fancyhdr}
\usepackage{xcolor}

% ── Page style ────────────────────────────────────────────────────────────────
\pagestyle{fancy}
\fancyhf{}
\fancyhead[L]{\textsc{STAT 452}}
\fancyhead[R]{\textsc{Lab 02}}
\fancyfoot[C]{\thepage}
\renewcommand{\headrulewidth}{0.4pt}

% ── Title formatting ─────────────────────────────────────────────────────────
\titleformat{\section}
  {\large\bfseries}{Exercise \thesection.}{0.5em}{}
\titleformat{\subsection}
  {\normalsize\bfseries}{}{0em}{}

% ── Document ──────────────────────────────────────────────────────────────────
\begin{document}

% ── Title block ───────────────────────────────────────────────────────────────
\begin{center}
  {\LARGE\bfseries Applied Statistics for Engineers and Scientists~II}\\[6pt]
  {\Large Lab 02 --- Practice Problems}\\[4pt]
  {\large Polynomial Regression, Transformations, Ridge, and Lasso}\\[8pt]
  {\normalsize Author: Ho Huu Binh \hspace{2em} Year: 2026}
\end{center}

\bigskip
\hrule
\bigskip

% ══════════════════════════════════════════════════════════════════════════════
\section{Transformation to linearize an exponential relationship}
% ══════════════════════════════════════════════════════════════════════════════

A laboratory records the number of surviving bacteria after radiation exposure:
$$
\begin{array}{c|cccccc}
\text{time }t & 0 & 1 & 2 & 3 & 4 & 5\\
\hline
\text{count }N_t & 100 & 80 & 65 & 51 & 41 & 33
\end{array}
$$
Assume an exponential mean law
$$
  N_t = N_0 e^{\beta t},
  \qquad N_0 > 0.
$$

\begin{enumerate}[label=(\alph*)]
  \item Linearize the model by applying an appropriate transformation.
        Write down the resulting linear model and identify the new intercept
        in terms of the original parameters.

  \item Fit the transformed regression in R using ordinary least squares.
        Report the estimated slope $\widehat\beta$ and the estimated initial
        count $e^{\widehat\alpha}$.

  \item Interpret the slope on the original (count) scale.
        What is the estimated multiplicative change in expected count for one
        additional time unit? What is the estimated proportional decrease per
        time unit?

  \item Check the residuals on the log scale.
        Explain why residual diagnostics should be performed on the
        transformed scale rather than the original count scale.
\end{enumerate}

\bigskip

% ══════════════════════════════════════════════════════════════════════════════
\section{Ridge shrinkage}
% ══════════════════════════════════════════════════════════════════════════════

Suppose a centered standardized design satisfies
$$
  X^\top X =
  \begin{pmatrix}
  10 & 8 \\
  8  & 10
  \end{pmatrix},
  \qquad
  X^\top y =
  \begin{pmatrix}
  9 \\
  9
  \end{pmatrix}.
$$

\begin{enumerate}[label=(\alph*)]
  \item Compute the ordinary least-squares estimate
        $\widehat\beta_{\mathrm{OLS}} = (X^\top X)^{-1} X^\top y$.

  \item Compute the ridge estimate for $\lambda = 2$:
        $$
          \widehat\beta_{\lambda} = (X^\top X + \lambda I_2)^{-1} X^\top y.
        $$

  \item Compute the eigenvalues of $X^\top X$ and of
        $X^\top X + 2I_2$.

  \item Compute the condition number
        $\kappa_2(A) = \lambda_{\max}(A) / \lambda_{\min}(A)$
        for both matrices. Explain why adding $2I_2$ improves numerical
        stability.
\end{enumerate}

\bigskip

% ══════════════════════════════════════════════════════════════════════════════
\section{Lasso soft-thresholding}
% ══════════════════════════════════════════════════════════════════════════════

Assume
$$
  X^\top X = I_4,
  \qquad
  X^\top y = (3,\; 0.7,\; -1.2,\; -4)^\top.
$$

\begin{enumerate}[label=(\alph*)]
  \item Using the soft-thresholding formula
        $$
          \widehat b_j
          = \operatorname{sign}(z_j)\bigl(|z_j| - \lambda\bigr)_+,
          \qquad
          z = X^\top y,
        $$
        compute the lasso solution $\widehat b_{\lambda}$ for $\lambda = 1$.
        Which variables are selected (i.e.\ have nonzero coefficients)?

  \item Repeat the computation for $\lambda = 2$.
        Which variables are selected now?

  \item Explain why increasing $\lambda$ makes the selected model smaller
        (i.e.\ fewer variables are retained).
\end{enumerate}

\end{document}
